% Compiled main-text proof export for TEXTBOOK-PART-IV-CHAPTER-16-INSTANCE-DEPENDENT-LOWER-BOUNDS-SPINE
\section{ABRL: A Target-Faithful Autoformalization Harness and Lean 4 Library for Bandit and Reinforcement Learning Theory}

\paragraph{Status.}
Definition~16.1, the source mean-to-gap producer, Theorem~16.2,
Lemma~16.3, and unit-variance Gaussian Theorem~16.4 compile.

\paragraph{Lean declarations.}
The source bridge is implemented by
\texttt{FiniteMeanBanditEnvironment},
\texttt{oneArmMeanChange\_produces\_gap\_contract}, and
\texttt{oneArmMeanIncrease\_sub\_gap\_eq\_changedMargin}.
The exact finite-time terminal is
\texttt{expectedPullCount\_ge\_log\_regret\_changeOfMeasure}.
The Gaussian specialization uses
\texttt{UnitVarianceGaussianBanditEnvironment},
\texttt{chapter16GaussianChangedEnvironment},
\texttt{chapter16GaussianChangedEnvironment\_armKL},
\texttt{InChapter16GaussianLocalClass},
\texttt{gaussianExpectedPullCount\_ge\_finiteTimeInstanceDependent},
\texttt{chapter16Gaussian\_finiteTime\_log\_identity}, and
\texttt{gaussianExpectedRegret\_ge\_finiteTimeInstanceDependent}.

\paragraph{Proof map.}
Tie every arm mean to the integral of its probability law. Unchanged arm laws
then have unchanged finite means. If arm $i$ is suboptimal in $\nu$, only its
law changes, and it is uniquely optimal in $\nu'$, the producer proves both
$\Delta_i(\nu)>0$ and
$\lambda-\Delta_i(\nu)>0$, together with the comparison to every unchanged-arm
gap in $\nu'$. The existing history-KL identity and the two exact
majority-event regret charges feed Bretagnolle--Huber, retaining the factor
$1/4$. Taking logarithms gives
\[
  \mathbb E_\nu[T_i(n)] \ge
  \frac{\log(\min\{\lambda-\Delta_i(\nu),\Delta_i(\nu)\}/4)
    +\log n-\log(R_n(\nu)+R_n(\nu'))}
       {D(P_i,P_i')}.
\]
The KL direction and expected-pull law are the source ones. Zero KL contradicts
the distinct certified means, while finite-positive and infinite KL are
handled as separate branches.

For Theorem~16.4, shift one positive-gap unit-Gaussian mean by
$(1+\epsilon)\Delta_i$. The changed vector lies in the coordinatewise local
class and its arm KL is $((1+\epsilon)\Delta_i)^2/2$. Lemma~16.3, the two
$Cn^p$ bounds, the exact logarithmic normalization, positive parts, and the
gap-times-pull decomposition yield
\[
 R_n(\nu) \ge \frac{2}{(1+\epsilon)^2}
 \sum_{i:\Delta_i>0}
 \left(\frac{(1-p)\log n+\log(\epsilon\Delta_i/(8C))}{\Delta_i}\right)^+.
\]
All source quantifiers on the horizon set, $C$, $p$, $\epsilon$, and the local
Gaussian class are retained.

\paragraph{Theorem 16.2.}
Replace one arm by any admissible strictly better finite-mean probability
law. Product-class membership is preserved. Consistency and Lemma~16.3
give the per-alternative reciprocal-KL liminf bound for the exact $n$-pull
regret and count sequences. In extended nonnegative reals,
$(\inf S)^{-1}=\sup_{a\in S}a^{-1}$ aggregates the costs, including empty,
zero-infimum, and infinite-cost branches. Finite-count Fatou and the exact
gap-times-pull decomposition yield
\[
 \liminf_{n\to\infty}\frac{R_n}{\log n}
 \ge \sum_{i:\Delta_i>0}\frac{\Delta_i}{d_{\inf}(P_i,\mu^*,\mathcal M_i)}.
\]
The compiled terminal is
\texttt{consistentPolicy\_liminf\_expectedRegret\_div\_log\_ge}.
